\section{Genuine minors and residual evidence capacity}
\label{sec:matroid-minors}

Gate~3 could express contraction only as a background-closure probe because the
lightweight \texttt{MatroidClosure} interface had no explicit ground set. Gate~4
removes that limitation. The formal layer now records a ground predicate together
with closure, ground containment, extensivity, monotonicity, idempotence, and
exchange, plus a finite list certifying that every ground atom occurs in a finite
support. The implementation is deliberately local to the repository and does not
add Mathlib as a dependency.

\begin{definition}[Explicit-ground deletion and contraction]
Let \(M\) be a ground-set matroid closure on \(E\), with ground set
\(E(M)\), and let \(D,C\subseteq E(M)\). Gate~4 defines
\[
  E(M\setminus D)=E(M)\setminus D,
  \qquad
  \operatorname{cl}_{M\setminus D}(A)
  =\operatorname{cl}_M(A)\setminus D,
\]
and
\[
  E(M/C)=E(M)\setminus C,
  \qquad
  \operatorname{cl}_{M/C}(A)
  =\operatorname{cl}_M(A\cup C)\setminus C.
\]
These are the standard closure formulas for deletion and contraction.
\end{definition}

\begin{theorem}[Minor closure preservation]
Both constructions again satisfy the explicit-ground closure axioms, including
Mac Lane--Steinitz exchange. The finite-ground wrappers retain a finite support
certificate for the smaller ground set.
\textup{\textsc{Status: Lean-proved in
\texttt{PEL4/MatroidEvidenceMinors.lean}; both constructors are axiom-free.}}
\end{theorem}

The explicit carrier makes the distinction between the two minor operations
mathematically literal. Deletion removes atoms from the problem. Contraction first
allows the contracted evidence to participate in generation and then removes it
from the residual carrier. It therefore represents a quotient by a fixed structural
direction, not the same operation as merely forgetting a source.

\subsection{Deletion: redundancy versus dispensability}

For the finite two-channel 4-PEL matroid, let \(x\) and \(y\) be distinct
same-polarity atoms in the ground set. Since \(\{x,y\}\) is a circuit, deleting
\(x\) leaves \(y\) as a representative of the same channel.

\begin{proposition}[Parallel deletion preserves the coarse state]
If \(x\neq y\), \(\chi(x)=\chi(y)\), and \(y\in E(M)\), then deleting \(x\)
preserves the presence or absence of both coarse support channels. Consequently,
for every FDE value \(v\),
\[
  \operatorname{RealizesFDE}(E(M)\setminus\{x\},v)
  \quad\Longleftrightarrow\quad
  \operatorname{RealizesFDE}(E(M),v).
\]
\textup{\textsc{Status: Lean-proved. The direct support-profile theorem is
constructive; the packaged FDE equivalence uses only the repository's allowed
standard Lean axioms.}}
\end{proposition}

This sharpens the Gate-3 notion of redundancy. Circuit membership implies that a
particular atom can be removed without changing the coarse observation only while
a parallel representative remains. Redundant relative to closure is therefore not
the same as globally dispensable under arbitrary subsequent deletions.

The complementary case is exact as well. Call \(c\) a unique channel
representative when it is the only ground atom of polarity \(\chi(c)\).

\begin{proposition}[Unique-representative deletion profile]
Deleting a unique channel representative removes exactly its own support channel;
every different polarity channel is preserved. In the bilateral four-valued
projection, the only possible one-atom information losses are therefore
\[
  \Tval\longrightarrow\Nval,
  \qquad
  \Fval\longrightarrow\Nval,
  \qquad
  \Bval\longrightarrow\Tval,
  \qquad
  \Bval\longrightarrow\Fval,
\]
with the direction determined by the deleted polarity and by which opposite
channel was present beforehand.
\textup{\textsc{Status: the exact channel-profile statement is Lean-proved and
axiom-free; the displayed four-value arrows are its immediate bilateral
corollaries.}}
\end{proposition}

Deletion thus exposes a useful distinction:
\[
  \text{coarse redundancy}
  \;\neq\;
  \text{unconditional removability}.
\]
An atom in a circuit is removable relative to another representative, but the last
representative of a channel is structurally pivotal for the FDE projection.

\subsection{Contraction: quotienting an independent direction}

The canonical finite channel matroid has no loops before taking minors. For a
single contracted atom \(c\), the general contraction identity gives
\[
  e\text{ is a loop in }M/c
  \quad\Longleftrightarrow\quad
  e\in E(M),\ e\neq c,
  \text{ and }e\in\operatorname{cl}_M(\{c\}).
\]
For the two-channel closure this becomes a complete polarity classification.

\begin{theorem}[Contraction loop profile]
For every remaining atom \(e\),
\[
  e\text{ is a loop in }M/c
  \quad\Longleftrightarrow\quad
  e\in E(M),\quad e\neq c,\quad \chi(e)=\chi(c).
\]
Thus contracting one positive atom turns every remaining positive parallel atom
into a loop, while no negative atom becomes a loop; conversely for negative
contraction.
\textup{\textsc{Status: Lean-proved and axiom-free in
\texttt{PEL4/MatroidEvidenceMinors.lean}.}}
\end{theorem}

This motivates a second support notion. Ordinary coarse support asks only whether
a polarity-labelled atom remains in the ground set. Define \emph{nonloop support}
to require a remaining ground atom of that polarity that is not a loop. Before
minor operations these notions coincide because the canonical channel matroid is
loopless.

\begin{proposition}[Residual channel capacity under contraction]
After contracting \(c\):
\begin{enumerate}[label=(\roman*)]
  \item the channel \(\chi(c)\) has no nonloop support;
  \item every channel different from \(\chi(c)\) has nonloop support after
  contraction exactly when it had ordinary ground support before contraction.
\end{enumerate}
Hence contraction quotients out exactly one coarse independent direction and leaves
the opposite independent direction active when present.
\textup{\textsc{Status: Lean-proved and axiom-free in
\texttt{PEL4/MatroidEvidenceMinorSemantics.lean}.}}
\end{proposition}

The distinction between ground support and nonloop support matters. A remaining
same-polarity atom may still be physically present in the contracted carrier while
contributing zero residual rank because it has become a loop. Consequently the
Gate-1 channel-rank reading as ``number of present support bits'' is exact for the
original loopless channel matroid but is not automatically invariant under minors.
After contraction, polarity occupancy and residual independent-information capacity
must be tracked separately.

\paragraph{Epistemic boundary.}
Genuine matroid contraction is a structural quotient. It is not, by itself, AGM
belief contraction, Bayesian conditioning, public announcement, source acceptance,
or a memory-bearing epistemic update. In particular, the contracted atom is absent
from the minor ground set. If an epistemic model is meant to remember that atom as
accepted background evidence, that memory must be represented by an additional
state component rather than silently identified with the matroid minor. Gate~4
therefore replaces the Gate-3 metaphor with a formal separation:
\[
  \text{source deletion}
  \neq
  \text{structural contraction}
  \neq
  \text{remembered background acceptance}.
\]

\paragraph{Next structural question.}
The finite-ground minor layer now makes it possible to ask whether the six-cell
reliability refinement admits a fine-grained matroid whose deletion and contraction
project homomorphically to the two-channel minors. That question would test whether
source reliability merely refines parallel classes or destroys the coarse minor
correspondence altogether.
